% TOP_PROBLEM: {"id": 30007428, "title": "Avoiding every geometric affine sequence in a large-measure set", "category_id": 8, "category": {"id": 8, "name": "analysis", "display_name": "Analysis", "description": "Limits, continuity, calculus, and function theory.", "slug": "analysis", "order_index": 8, "created_at": "2026-07-31T15:26:25.671Z"}, "status": "open"}
% FIRST_POSTED: 2026-10-06
% ATTEMPT_STATUS: solved
% OPENAI_CLAIM: {"schema_version":1,"source_repo":"https://github.com/openai/math","source_commit":"adc7f1241b42e322a6451854ab7e4b4c146bf78a","release_date":"2026-10-06","match":"full","resolution":"proved","independently_reviewed":false,"second_pass":"agreed","families":[{"family":"084","title":"The geometric case of the Erdős similarity conjecture","subject":"Real and complex analysis","manuscripts":[{"title":"The geometric case of the Erdős similarity conjecture","pdf_path":"preprints/The-geometric-case-of-the-Erdos-similarity-conjecture-October-5-2026/geometric-erdos-similarity.pdf","readme_path":"preprints/The-geometric-case-of-the-Erdos-similarity-conjecture-October-5-2026/README.md","date":"2026-10-05","theorem_ref":"thm:main (sections/01-introduction.tex:17)","statement_sources":[{"path":"preprints/The-geometric-case-of-the-Erdos-similarity-conjecture-October-5-2026/build/sections/01-introduction.tex","line":17,"label":"thm:main"}]}],"lean":{"doc_path":"lean/docs/084.md","comparators":[],"declarations":[],"covers_main_theorem":false},"reasoning_trace":null}]}
\documentclass[11pt]{article}
\usepackage[margin=25mm]{geometry}
\usepackage{fontspec}
\setmainfont{Latin Modern Roman}
\usepackage{amsmath,amssymb,mathtools}
\usepackage{unicode-math}
\setmathfont{Latin Modern Math}
\usepackage{enumitem}
\usepackage{xurl,hyperref,bookmark}
\hypersetup{colorlinks=true,urlcolor=blue}
\setcounter{secnumdepth}{0}
\setlength{\parindent}{0pt}
\setlength{\parskip}{5pt}
\setlength{\emergencystretch}{3em}
\begin{document}
\section*{Avoiding every geometric affine sequence in a large-measure set}

\subsection{OpenAI's claim}
\textit{Avoiding a fixed geometric progression}: For each fixed ratio $q$ strictly between zero and one, the manuscript claims almost-full-measure compact sets in the unit interval avoiding every affine copy of $(q^n)$. The chosen set is permitted to depend on the ratio.

\subsection{Mathematical statement of the claimed result}
\paragraph{Mathematical claim 1.} For any fixed $0<q<1$ and $0<\eta<1$ the claimed construction produces compact $E_{q,\eta}\subset[0,1]$ with $m(E_{q,\eta})>1-\eta$ and \[\forall x\in\mathbb R\ \forall s\in\mathbb R\setminus\{0\}\ \exists n\ge1:\quad x+sq^n\notin E_{q,\eta}.\] The avoiding set may depend on both $q$ and $\eta$.

\subsection{Outline of the argument}
\textit{Avoiding a fixed geometric progression}: The construction adapts the ordered-tree random hitting method, using nested grids with resolutions matched to the chosen geometric ratio. Gaps between edge index intervals preserve previous routing decisions and expose independent current trials. The number of scale representatives is controlled locally within each subtree, and a final neighborhood repair deals with centers missed by the random construction.

\subsection{Scope relative to this problem}
For each fixed $q\in(0,1)$ and each positive measure tolerance, the theorem constructs a compact subset of $[0,1]$ avoiding every affine copy of $\{q^n:n\ge1\}$ with nonzero dilation of either sign. Its measure exceeds one minus the tolerance. This is exactly the geometric-progression assertion recorded here, with the avoiding set allowed to depend on both parameters.

\subsection{Formal verification}
Family 084 describes Lean coverage for avoidance for the dyadic ratio; it does not assert avoidance for an arbitrary geometric ratio.

\textit{Avoiding a fixed geometric progression}: No manuscript-specific Comparator is listed in the pinned inventory.

The documented selection does not cover the entire manuscript assertion relevant here.

This site has not run the Lean development. These links describe OpenAI’s selected formal statements and should be read with their stated scope.

\subsection{Status caveat}
OpenAI reports that an internal model produced the manuscripts. Their arguments have not been independently verified as proofs by this site. Any portion outside the documented formal scope may contain mistakes, and the catalogue label reports the announced mathematical claim.

\subsection{Sources}
\begin{itemize}
\item OpenAI, \emph{The geometric case of the Erdős similarity conjecture}, 2026-10-05; citation key \texttt{OAI:The-geometric-case-of-the-Erdos-similarity-conjecture-October-5-2026}. \href{https://github.com/openai/math/blob/adc7f1241b42e322a6451854ab7e4b4c146bf78a/preprints/The-geometric-case-of-the-Erdos-similarity-conjecture-October-5-2026/geometric-erdos-similarity.pdf}{PDF (GitHub)}; \href{https://github.com/openai/math/raw/adc7f1241b42e322a6451854ab7e4b4c146bf78a/preprints/The-geometric-case-of-the-Erdos-similarity-conjecture-October-5-2026/geometric-erdos-similarity.pdf}{PDF (direct)}; \href{https://github.com/openai/math/blob/adc7f1241b42e322a6451854ab7e4b4c146bf78a/preprints/The-geometric-case-of-the-Erdos-similarity-conjecture-October-5-2026/README.md}{Manuscript page and citation}; \href{https://github.com/openai/math/blob/main/preprints/The-geometric-case-of-the-Erdos-similarity-conjecture-October-5-2026/geometric-erdos-similarity.pdf}{Latest version}; Statement source: \href{https://github.com/openai/math/blob/adc7f1241b42e322a6451854ab7e4b4c146bf78a/preprints/The-geometric-case-of-the-Erdos-similarity-conjecture-October-5-2026/build/sections/01-introduction.tex\#L17}{thm:main}.
\item Family 084: \href{https://github.com/openai/math/blob/adc7f1241b42e322a6451854ab7e4b4c146bf78a/lean/docs/084.md}{Lean scope}.
\item \href{https://github.com/openai/math/blob/adc7f1241b42e322a6451854ab7e4b4c146bf78a/CONTENTS.md}{OpenAI catalogue entry}.
\end{itemize}
\end{document}
